\documentclass[conference,a4paper]{IEEEtran}

\usepackage[T1]{fontenc}
\usepackage[utf8]{inputenc}
\usepackage[brazil]{babel}
\usepackage{amsmath,amssymb}
\usepackage{booktabs}
\usepackage{graphicx}
\usepackage{microtype}
\usepackage{url}
\usepackage{placeins}

\graphicspath{{../resultados/}}

\title{Avaliação de Políticas de Balanceamento de Carga em um Sistema Distribuído}

\author{%
  \IEEEauthorblockN{Caique Pinheiro Andrade e Luiz Henrique Marques Gonçalves}
  \IEEEauthorblockA{Instituto de Computação -- Universidade Estadual de Campinas\\
  MC714 -- Sistemas Distribuídos -- 2º semestre de 2026\\
  \texttt{c204677@dac.unicamp.br, l183218@dac.unicamp.br}}
}

\begin{document}
\maketitle

\begin{abstract}
Este trabalho avalia, por modelo analítico e simulação, um balanceador que
distribui chegadas Poisson entre três servidores homogêneos, cada um com uma
fila FCFS e serviço exponencial de taxa $\mu=1$. Comparam-se as políticas
aleatória, Round Robin e fila mais curta. Na aleatória, a decomposição do
processo de Poisson produz três filas M/M/1 independentes de taxa
$\lambda/3$. Nas 15 configurações estáveis (10 réplicas cada), vazão e
utilização diferiram do previsto em no máximo 0,45\% e 0,66\%, e a Lei de
Little teve erro máximo de 0,13\%. A fila mais curta reduziu o tempo médio de
resposta em até 58,3\%. Com $\lambda=3{,}3>3\mu$, a vazão saturou em
3 req./u.t.\ e $N(t)$ cresceu aproximadamente como $(\lambda-3\mu)t$.
\end{abstract}

\section{Introdução}

Em sistemas distribuídos, o balanceamento distribui requisições para evitar
a sobrecarga de alguns servidores e a ociosidade de outros. Este trabalho
implementa o sistema proposto, modela analiticamente a escolha aleatória e a
compara, por simulação, com Round Robin e fila mais curta. A Seção II apresenta o sistema e as
políticas; a Seção III desenvolve os itens analíticos (a)--(c); as Seções IV
e V descrevem o simulador e o desenho experimental; a Seção VI discute os
resultados dos itens (d)--(f); e a Seção VII apresenta as conclusões.

\section{Sistema e políticas}

\subsection{Modelo do sistema}

O sistema tem dois componentes. O balanceador recebe todas as requisições e
é instantâneo: decide o destino e encaminha sem consumir tempo. Os três
servidores são homogêneos; cada um tem uma única
thread de execução, atende uma requisição por vez e mantém as demais em uma
fila FCFS (\emph{First Come, First Served}) ilimitada.

As chegadas ao balanceador são Poisson de taxa $\lambda$; portanto, os
tempos entre chegadas são exponenciais independentes de média $1/\lambda$.
O serviço é exponencial de média $1/\mu$, com $\mu=1$ req./u.t.\ e
$E[S]=1$ u.t. A intensidade por servidor é
$\rho=\lambda/(3\mu)$, e a estabilidade exige $\rho<1$, ou
$\lambda<3\mu$.

\subsection{Políticas de roteamento}

Foram implementadas três políticas de balanceamento, selecionáveis pela
linha de comando do simulador: \textbf{aleatória (A)}, que sorteia o
servidor de cada requisição com probabilidade $1/3$, de forma
independente; \textbf{Round Robin (RR)}, que encaminha ciclicamente
($1,2,3,1,2,3,\dots$) sem consultar o estado; e \textbf{fila mais curta
(FC)}, que consulta a ocupação de cada servidor (em execução mais em
espera) e envia à menor, sorteando entre os empatados. Logs periódicos
registram a ocupação das filas e quantas requisições cada servidor recebeu,
permitindo acompanhar a dinâmica da distribuição de carga.

\section{Modelo analítico}

Esta seção desenvolve os itens (a)--(c) para a política aleatória; RR e FC
são comparadas ao modelo por simulação.

\subsection{(a) Decomposição das chegadas}

Seja $N(t)$ o número de chegadas ao balanceador até o instante $t$. Por
hipótese, $N(t)$ é um processo de Poisson de taxa $\lambda$. Na política
aleatória, cada chegada recebe, independentemente, a marca
$i\in\{1,2,3\}$ do servidor escolhido, com probabilidade $p_i=1/3$. Seja
$N_i(t)$ a contagem das chegadas que receberam a marca $i$.

O teorema da decomposição (\emph{splitting}) afirma que a marcação
independente de um processo de Poisson produz um processo de Poisson para
cada marca. A taxa do subprocesso $i$ é a taxa original multiplicada pela
probabilidade da marca. De fato, em um intervalo de duração $\Delta t$
ocorrem, em média, $\lambda\Delta t$ chegadas; a fração média destinada ao
servidor $i$ é $p_i$, resultando em $p_i\lambda\Delta t$ chegadas. Como
$p_i=1/3$,
\begin{equation}
 \lambda_i=p_i\lambda=\frac{\lambda}{3},\qquad i=1,2,3.
 \label{eq:splitting}
\end{equation}
O cálculo da média determina a taxa, enquanto o teorema garante que cada
$N_i(t)$ continua sendo Poisson e que os três subprocessos são mutuamente
independentes. Essa independência decorre da marcação: o servidor sorteado
para uma chegada não influencia o sorteio das demais. Portanto, sob a
política aleatória, cada servidor recebe um processo de Poisson independente
de taxa $\lambda/3$.

\subsection{(b) Filas M/M/1, métricas e estabilidade}

O servidor $i$ recebe chegadas Poisson de taxa $\lambda/3$, atende uma
requisição por vez com tempo exponencial de taxa $\mu$ e mantém disciplina
FCFS: essa é exatamente uma fila M/M/1. Seja $N_i$ o número de
requisições no servidor $i$ (em execução e em espera) no regime estacionário
e seja $p_k=P[N_i=k]$ a probabilidade de haver exatamente $k$ requisições
nele, para $k=0,1,2,\dots$. Retomando a distribuição estacionária fornecida,
\begin{equation}
 p_k=(1-\rho)\rho^k,\qquad
 \rho=\frac{\lambda_i}{\mu}=\frac{\lambda}{3\mu}.
 \label{eq:pk}
\end{equation}
Para que essa seja uma distribuição de probabilidade, a série
$\sum_{k\ge0}p_k$ deve convergir; isso ocorre somente se $\rho<1$. Portanto,
a condição de estabilidade é $\lambda<3\mu$. Sob essa condição, as métricas
pedidas seguem diretamente da distribuição em (\ref{eq:pk}). Para o estado
vazio, (\ref{eq:pk}) dá $p_0=(1-\rho)\rho^0=1-\rho$; logo, a probabilidade
de o servidor estar ocupado é $U_i=P(N_i>0)=1-p_0=\rho$. Para o número médio de requisições, substituímos (\ref{eq:pk}) na definição
de esperança e usamos $\sum_{k\ge0}k\rho^k=\rho/(1-\rho)^2$:
\begin{align}
 E[N_i]
 &=\sum_{k\ge0}k\,p_k \\[-1mm]
 &= (1-\rho)\sum_{k\ge0}k\rho^k \\[-1mm]
 &= (1-\rho)\frac{\rho}{(1-\rho)^2}
  =\frac{\rho}{1-\rho}.
 \label{eq:eni}
\end{align}
Para cada servidor $i$, a vazão é limitada simultaneamente pela chegada e pela capacidade
do servidor: $X_i=\min(\lambda/3,\mu)=\lambda/3$, pois $\lambda<3\mu$. Essa é a taxa de saída
usada na Lei de Little para o servidor $i$; portanto,
$E[N_i]=\lambda_iE[R]$. Assim,
\begin{equation}
 E[R]=\frac{E[N_i]}{\lambda_i}
 =\frac{\rho/(1-\rho)}{\lambda_i}
 =\frac{1}{\mu-\lambda/3}.
 \label{eq:er}
\end{equation}
Como uma resposta é composta pela espera $T_Q$ mais o serviço $S$,
$R=T_Q+S$. Tomando esperanças e usando $E[S]=1/\mu$,
\begin{equation}
 E[T_Q]=E[R]-E[S]
 =\frac{1}{\mu-\lambda/3}-\frac{1}{\mu}
 =\frac{\rho}{\mu-\lambda/3}.
 \label{eq:etq}
\end{equation}
Como as filas são independentes e idênticas, $E[N]=3E[N_i]$. A
Tabela~\ref{tab:analitico} reúne essas previsões para os $\lambda$ do
experimento.

\begin{table}[!ht]
\caption{Previsões analíticas da escolha aleatória ($\mu=1$).}
\label{tab:analitico}
\centering
\setlength{\tabcolsep}{2.4pt}
\footnotesize
\begin{tabular}{@{}rrrrrrr@{}}
\toprule
$\lambda$ & $\rho$ & $U_i$ & $E[N_i]$ & $E[T_Q]$ & $E[R]$ & $X$ \\
\midrule
0,6 & 0,2 & 0,2 & 0,250 & 0,250 & 1,250 & 0,6 \\
1,2 & 0,4 & 0,4 & 0,667 & 0,667 & 1,667 & 1,2 \\
1,8 & 0,6 & 0,6 & 1,500 & 1,500 & 2,500 & 1,8 \\
2,4 & 0,8 & 0,8 & 4,000 & 4,000 & 5,000 & 2,4 \\
2,7 & 0,9 & 0,9 & 9,000 & 9,000 & 10,000 & 2,7 \\
\bottomrule
\end{tabular}
\end{table}
\FloatBarrier

\subsection{(c) Vazão e utilização valem para as três políticas}

O roteamento escolhe \emph{qual} servidor atenderá, mas não altera as
chegadas nem o trabalho de cada requisição. Em uma janela longa $T$, entram
em média $\lambda T$ requisições. No regime estável não há acúmulo; como não
há descartes, partem também $\lambda T$. Dividindo por $T$,
\begin{equation}
 X=\lambda
 \label{eq:vazao-conservacao}
\end{equation}
para as três políticas.

A Lei de Little aplicada ao subsistema das requisições \emph{em execução},
 que recebe todas elas, à taxa $X=\lambda$, e as retém por $E[S]=1/\mu$,
dá $\lambda/\mu$ como número médio de requisições em serviço. Como o
servidor $i$ está ocupado exatamente quando tem uma em execução, $U_i$ é a
probabilidade disso, e a soma das utilizações é esse mesmo número médio:
\begin{equation}
 U_1+U_2+U_3=\frac{\lambda}{\mu}.
 \label{eq:utilizacao-agregada}
\end{equation}
Os servidores são homogêneos e nenhuma política privilegia índices: A
sorteia uniformemente, RR mantém contagens que diferem em no máximo uma e FC
sorteia empates. Logo, $U_1=U_2=U_3$ e
\begin{equation}
 U_i=\frac{\lambda}{3\mu}=\rho.
 \label{eq:utilizacao-conservacao}
\end{equation}
Esses argumentos de conservação, que não exigem filas M/M/1, valem para A,
RR e FC; a Tabela~\ref{tab:simulacao} os confirma.

\section{Arquitetura do simulador}

\subsection{Simulação por eventos discretos}

Na simulação por eventos discretos, o relógio salta ao próximo evento, sem
passos fixos nem erro de discretização. As três threads são objetos lógicos
\texttt{Servidor}, não threads do sistema operacional; assim, sincronização
e escalonamento reais não interferem no modelo.

\subsection{Fila de eventos e tratamento}

Os eventos futuros ficam em um heap binário, uma fila de prioridade que
remove primeiro o evento de menor instante. Cada item é ordenado por
(instante, tipo, ordem de inserção). Há três tipos: \texttt{CHEGADA}, quando
uma requisição alcança o balanceador; \texttt{PARTIDA}, quando um serviço
termina; e \texttt{LOG}, que apenas registra as ocupações periodicamente. Em
empates temporais, a partida vem primeiro e libera o servidor antes da nova
alocação.

Cada servidor mantém uma \texttt{deque} de pares (instante de chegada,
tempo de serviço). O primeiro item é a requisição em execução; os demais
aguardam atrás dele em ordem FCFS. Em uma \texttt{CHEGADA}$(t)$, o
balanceador lê as ocupações, aplica a política e insere a nova requisição no
fim da fila. Se ela estava vazia, inicia o serviço e agenda a partida para
$t+\text{serviço}$. Em uma \texttt{PARTIDA}$(t)$, remove a requisição em
execução e registra $R=t-t_{\mathrm{chegada}}$; se houver outra aguardando,
ela inicia imediatamente e sua partida é agendada. Assim, cada servidor
executa no máximo uma requisição e preserva a ordem FCFS.

\subsection{Coleta de métricas}

As métricas cobrem $[500,5000]$, uma janela de 4500 u.t.; as primeiras
500 u.t.\ são aquecimento (\emph{warm-up}) e removem o transiente inicial.
Antes de cada evento, o simulador integra o estado vigente: se $N(t)=n$ por
$\Delta t$, soma $n\Delta t$ à área de $N$; se o servidor $i$ está ocupado,
soma $\Delta t$ ao seu tempo ocupado. Intervalos que cruzam o instante 500
são recortados, contabilizando-se apenas a parcela posterior ao aquecimento.
Assim,
$E[N]=\int_{500}^{5000}N(t)\,dt/4500$ e
$U_i=\text{tempo ocupado}_i/4500$. Por fim,
$X=\text{partidas}/4500$, e $E[R]$ é a média das respostas das requisições
que partem na janela. Desse modo, $E[N]$ é uma média ao longo do tempo, e
não uma média das ocupações observadas apenas nas chegadas.

\subsection{Reprodutibilidade e validação}

Três geradores separados produzem chegadas, serviços e decisões de política.
Essa separação impede que os sorteios de uma política alterem a sequência de
chegadas ou serviços. O serviço é sorteado na chegada e acompanha a
requisição. Numa réplica, as políticas
recebem os mesmos instantes e demandas, mudando apenas o roteamento. Esses
números aleatórios comuns tornam a comparação pareada e reduzem sua
variância. Os testes automatizados verificam RR, empates da FC,
conservação, reprodutibilidade, $E[R]\to E[S]$ em carga baixa, Lei de Little
e os regimes estável e instável. Instruções de execução estão no README.

\section{Desenho experimental}

As 15 configurações combinam
$\lambda\in\{0{,}6;1{,}2;1{,}8;2{,}4;2{,}7\}$ com as três políticas. Cada
uma teve 10 sementes, 5000 u.t.\ de duração e 500 u.t.\ de aquecimento. Para uma
métrica, $\bar{x}$ é a média e $s$ o desvio-padrão das réplicas. Como há
apenas 10 amostras e o desvio populacional é desconhecido, usa-se a
$t$ de Student com $n-1=9$ graus de liberdade. Suas caudas mais largas
incorporam a incerteza de estimar o desvio por $s$. O IC
bilateral de 95\% é
\begin{equation}
 \bar{x}\ \pm\ t_{0{,}975;9}\,\frac{s}{\sqrt{10}},
 \qquad t_{0{,}975;9}=2{,}262,
 \label{eq:ic}
\end{equation}
onde $t_{0{,}975;9}=2{,}262$ delimita os 95\% centrais da distribuição
$t_9$. As barras representam a incerteza da média, não a variação
individual. O cenário $\lambda=3{,}3$ parte vazio, sem aquecimento, e registra
$N(t)$ para comparação com a aproximação fluida.

\section{Resultados e discussão}

\subsection{Vazão e utilização (item c)}

A Tabela~\ref{tab:simulacao} traz média $\pm$ semiamplitude do IC de todas
as configurações. A vazão simulada diferiu de $\lambda$ em no máximo
0,45\%, e a utilização média diferiu de $\rho$ em no máximo 0,66\%; as
utilizações individuais são equilibradas entre servidores. Isso confirma o
argumento de conservação da Seção III-C: $X$ e $U_i$ independem da política
de roteamento.

\begin{table*}[t]
\caption{Média $\pm$ semiamplitude do IC de 95\% em dez réplicas. A, RR e FC
denotam Aleatória, Round Robin e Fila Mais Curta; $X E[R]$ é calculado pelo
produto das médias.}
\label{tab:simulacao}
\centering
\scriptsize
\renewcommand{\arraystretch}{0.95}
\setlength{\tabcolsep}{2.2pt}
\begin{tabular}{@{}rrlllllll@{}}
\toprule
\multicolumn{1}{c}{$\lambda$} & \multicolumn{1}{c}{Pol.} &
\multicolumn{1}{c}{$X$} & \multicolumn{1}{c}{$E[R]$} &
\multicolumn{1}{c}{$E[N]$} & \multicolumn{1}{c}{$XE[R]$} &
\multicolumn{1}{c}{$U_1$} & \multicolumn{1}{c}{$U_2$} &
\multicolumn{1}{c}{$U_3$} \\
\midrule
0,6 & A  & $0,598\pm0,006$ & $1,237\pm0,025$ & $0,741\pm0,017$ & $0,741$ & $0,196\pm0,007$ & $0,200\pm0,008$ & $0,200\pm0,007$ \\
    & RR & $0,598\pm0,006$ & $1,058\pm0,017$ & $0,633\pm0,011$ & $0,633$ & $0,197\pm0,005$ & $0,201\pm0,006$ & $0,198\pm0,005$ \\
    & FC & $0,598\pm0,006$ & $1,019\pm0,016$ & $0,610\pm0,011$ & $0,610$ & $0,196\pm0,003$ & $0,200\pm0,008$ & $0,200\pm0,006$ \\
1,2 & A  & $1,205\pm0,012$ & $1,678\pm0,030$ & $2,022\pm0,048$ & $2,023$ & $0,402\pm0,009$ & $0,397\pm0,009$ & $0,405\pm0,008$ \\
    & RR & $1,205\pm0,012$ & $1,292\pm0,024$ & $1,558\pm0,033$ & $1,558$ & $0,402\pm0,006$ & $0,405\pm0,008$ & $0,397\pm0,005$ \\
    & FC & $1,205\pm0,012$ & $1,135\pm0,017$ & $1,368\pm0,025$ & $1,368$ & $0,403\pm0,003$ & $0,402\pm0,008$ & $0,399\pm0,007$ \\
1,8 & A  & $1,801\pm0,014$ & $2,517\pm0,109$ & $4,533\pm0,203$ & $4,533$ & $0,600\pm0,011$ & $0,605\pm0,010$ & $0,588\pm0,010$ \\
    & RR & $1,801\pm0,014$ & $1,796\pm0,043$ & $3,234\pm0,091$ & $3,235$ & $0,595\pm0,012$ & $0,601\pm0,007$ & $0,596\pm0,009$ \\
    & FC & $1,801\pm0,015$ & $1,404\pm0,027$ & $2,528\pm0,057$ & $2,529$ & $0,597\pm0,007$ & $0,600\pm0,007$ & $0,595\pm0,008$ \\
2,4 & A  & $2,408\pm0,011$ & $5,108\pm0,188$ & $12,314\pm0,482$ & $12,302$ & $0,811\pm0,010$ & $0,791\pm0,013$ & $0,809\pm0,010$ \\
    & RR & $2,409\pm0,011$ & $3,505\pm0,160$ & $8,449\pm0,407$ & $8,443$ & $0,803\pm0,008$ & $0,800\pm0,012$ & $0,807\pm0,011$ \\
    & FC & $2,409\pm0,012$ & $2,316\pm0,086$ & $5,582\pm0,220$ & $5,580$ & $0,803\pm0,007$ & $0,805\pm0,007$ & $0,803\pm0,007$ \\
2,7 & A  & $2,691\pm0,024$ & $9,110\pm0,706$ & $24,532\pm2,020$ & $24,517$ & $0,889\pm0,013$ & $0,888\pm0,016$ & $0,908\pm0,011$ \\
    & RR & $2,691\pm0,023$ & $6,496\pm0,498$ & $17,487\pm1,426$ & $17,480$ & $0,900\pm0,013$ & $0,893\pm0,008$ & $0,892\pm0,014$ \\
    & FC & $2,691\pm0,023$ & $3,799\pm0,280$ & $10,235\pm0,812$ & $10,222$ & $0,895\pm0,009$ & $0,894\pm0,009$ & $0,896\pm0,010$ \\
\bottomrule
\end{tabular}
\end{table*}

\subsection{Tempo de resposta e políticas (item d)}

A Fig.~\ref{fig:comparacoes} compara o $E[R]$ simulado das três políticas
com a curva M/M/1 da aleatória, dada por (\ref{eq:er}). Apenas ela recebe,
em cada servidor, chegadas Poisson independentes. No RR, as chegadas por
servidor têm intervalos Erlang-3; na FC, o destino depende do estado das três
filas, que ficam acopladas. Portanto, a curva é uma previsão exata somente
para a aleatória; RR e FC são comparadas a ela experimentalmente. Em todas
as cinco cargas, observou-se
\begin{equation}
 E[R]_{\mathrm{FC}}\le E[R]_{\mathrm{RR}}\le E[R]_{\mathrm{A}}
 \approx\frac{1}{\mu-\lambda/3}.
 \label{eq:ordem}
\end{equation}
A explicação está na forma como cada política distribui as chegadas. No RR,
cada servidor recebe uma a cada três chegadas globais. O intervalo entre suas
chegadas é, portanto, a soma de três exponenciais (Erlang-3), cujo
coeficiente de variação $1/\sqrt{3}$ é menor que o valor 1 da exponencial;
menos rajadas reduzem a espera. A FC usa
informação adicional: observa a ocupação atual e evita enviar uma requisição
para uma fila maior enquanto existe uma menor. Por isso, reduz diretamente o
desequilíbrio entre os servidores e obtém a menor espera.

Os ganhos de RR frente à aleatória, para
$\lambda=(0{,}6;1{,}2;1{,}8;2{,}4;2{,}7)$, foram
$(14{,}46;22{,}98;28{,}65;31{,}38;28{,}69)$\%; os de FC,
$(17{,}68;32{,}36;44{,}22;54{,}66;58{,}30)$\%. O ganho percentual de uma
política $p$ é definido por
$100\,(E[R]_{\mathrm{A}}-E[R]_p)/E[R]_{\mathrm{A}}$. Quanto maior a carga,
mais caro é enviar uma requisição a um servidor ocupado enquanto outro está
mais livre -- o erro que a FC evita; o RR ganha menos porque apenas
regulariza as chegadas, sem consultar o estado.

\begin{figure}[t]
\centering
\includegraphics[width=0.95\columnwidth]{er_vs_lambda.pdf}
\caption{Tempo médio de resposta $E[R]$ em função de $\lambda$. A curva
tracejada é a previsão M/M/1 da política aleatória; os pontos são médias de
dez réplicas e as barras, ICs de 95\%.}
\label{fig:comparacoes}
\end{figure}

A política aleatória aderiu ao modelo: até $\lambda=2{,}4$, o desvio do
$E[R]$ foi de no máximo 2,16\%, e o IC contém a curva em todos os quatro
pontos. Em $\lambda=2{,}7$, o valor simulado foi $9{,}110\pm0{,}706$ e o IC
$[8{,}40;9{,}82]$ não contém a previsão 10 (desvio de $-8{,}9$\%). 

O desvio em $\lambda=2{,}7$ é o viés de horizonte finito, não um erro do simulador. A execução parte com as filas vazias, muito abaixo do nível estacionário $E[N_i]=9$ correspondente a $\rho=0{,}9$; como a carga está próxima da saturação, as flutuações são grandes e a influência do estado inicial decai lentamente, de modo que as 500 u.t.\ de aquecimento não a eliminam e parte da janela ainda mede filas menores que as de equilíbrio. O padrão dos desvios confirma o diagnóstico: $E[N]$ e $E[R]$ caem na mesma proporção ($-9{,}1$\% e $-8{,}9$\%) enquanto $X$ permanece correto, que é exatamente o efeito de um transiente de fila, a vazão é fixada pela taxa de chegada, mas o nível das filas ainda não subiu. Alongando o horizonte, o viés desaparece: $E[R]=9{,}754\pm0{,}499$ com 20000 u.t., e $E[R]=9{,}997\pm0{,}321$ com $E[N]=27{,}01\pm0{,}89$ com 50000 u.t.\ e 5000 u.t.\ de aquecimento, contra os 10 e 27 previstos. Para cargas menores o nível estacionário é mais baixo e o viés é desprezível. 

\subsection{Lei de Little (item e)}

A Lei de Little, $E[N]=X\,E[R]$, vale para qualquer sistema estável. Ela é
verificada com estimadores independentes: $E[N]$ vem da integral temporal de
$N(t)$; $X$, da contagem de partidas; e $E[R]$, dos tempos individuais de
resposta. A Tabela~\ref{tab:simulacao} apresenta os dois lados da igualdade
para todas as configurações. As colunas $E[N]$ e $X E[R]$ são praticamente
idênticas, com erro relativo máximo de 0,13\%. Essa concordância confirma a
Lei de Little e também valida a integração temporal de $N(t)$ e a delimitação
da janela de medição.

\subsection{Regime instável (item f)}

Para $\lambda=3{,}3>3\mu=3$, a chegada excede a capacidade agregada e
$\rho=1{,}1>1$. A série geométrica de (\ref{eq:pk}) diverge: não há regime
estacionário, e $E[N]$, $E[T_Q]$ e $E[R]$ crescem sem limite. Portanto, as
fórmulas estacionárias da Seção~III-B não se aplicam. Depois que os três
servidores ficam ocupados, o sistema recebe a taxa $\lambda$ e escoa no
máximo $3\mu$; a aproximação fluida é, portanto,
\begin{equation}
 N(t)\approx N(0)+(\lambda-3\mu)t=N(0)+0{,}3t.
 \label{eq:fluida}
\end{equation}
Na simulação, as vazões de A, RR e FC foram 3,001, 3,004 e 3,007 req./u.t.:
a saída satura na capacidade, não na chegada. As inclinações médias de
$N(t)$ foram 0,331, 0,329 e 0,326 req./u.t., contra os 0,3 da reta fluida
nominal (Fig.~\ref{fig:instavel}). A diferença é da réplica, não do modelo:
este cenário é uma execução única, na qual chegaram 16663 requisições em
5000 u.t., ou seja $\hat{\lambda}=3{,}333$, uma flutuação de $1{,}3$
desvios-padrão em uma contagem de Poisson de média 16500 e desvio
$\sqrt{16500}\approx128$. Substituindo $\lambda$ por $\hat{\lambda}$ em
(\ref{eq:fluida}), a inclinação prevista passa a 0,333 e o acúmulo final a
1663 requisições, dos quais os três valores simulados (1656, 1643 e 1630)
distam no máximo 2,0\%. As três curvas quase se sobrepõem: com todos os
servidores permanentemente ocupados, o roteamento torna-se irrelevante.

\begin{figure}[t]
\centering
\includegraphics[width=0.95\columnwidth]{instavel_Nt.pdf}
\caption{Número de requisições no sistema, $N(t)$, no regime instável
$\lambda=3{,}3$, com a aproximação fluida $(\lambda-3\mu)t$ para
comparação.}
\label{fig:instavel}
\end{figure}

\noindent\begin{minipage}{\columnwidth}
\section{Conclusão}

O simulador de eventos discretos do balanceador foi validado de três formas:
as invariantes de conservação ($X$ e $U_i$ iguais para as três políticas,
desvios abaixo de 0,7\%), a adesão da política aleatória ao modelo M/M/1
(erro máximo de 2,2\% até $\lambda=2{,}4$; o desvio em $\rho=0{,}9$ é o
transiente de horizonte finito) e a Lei de Little (erro máximo de 0,13\%). Entre as políticas, a fila mais curta é a
melhor em todas as cargas -- até 58,3\% de redução no tempo médio de
resposta -- por usar a ocupação instantânea, enquanto o RR melhora apenas
regularizando as chegadas. Já o cenário
$\lambda=3{,}3>3\mu$ mostrou o limite de qualquer roteamento: a vazão
satura na capacidade agregada e as filas divergem linearmente. Rotear bem
reduz a espera, mas não substitui capacidade.

\section*{Divisão de trabalho}
\noindent\textit{Caique Pinheiro Andrade:} implementação do simulador e das
políticas, e execução da campanha experimental.
\textit{Luiz Henrique Marques Gonçalves:} derivação analítica, análise dos
resultados e redação do relatório.
\end{minipage}

\end{document}
